\documentclass[twocolumn]{aastex631}
\begin{document}
\title{A residual reionization ionizing-photon-budget shortfall for star-forming galaxies at z\~{}8 (lowxi systematic corner)}
\author{NebulaMind Lab (autonomous pipeline)}
\affiliation{NebulaMind Lab, \url{https://lab.nebulamind.net}}
\begin{abstract}
Reionization ionizing-photon-budget reconciliation at z\~{}8: star-forming galaxies require f\_esc=0.419 (+0.422/-0.214) to close the budget (Madau-Dickinson SFRD, log xi\_ion=25.5+/-0.15, clumping C=2-5, JWST-SFRD tail), versus indirect-proxy-inferred f\_esc=0.062 (+0.108/-0.039) from LzLCS O32/beta calibrations. Median delta(required-inferred)=+0.330 dex-frac (16-84\%: +0.100 to +0.752); 93\% of the systematic MC shows a shortfall. Conclusion: a genuine SHORTFALL remains, and the sign holds under both O32 and beta calibrations. The result is bounded by the xi\_ion x clumping x proxy-calibration systematic, not by statistics. Generated autonomously from public data (jwst) via the NebulaMind Lab runner: a bounded, reproducible, descriptive study.
\end{abstract}
\section{Introduction}
Introduction:
Recent studies have highlighted a potential crisis in understanding the reionization process, with concerns that current models may not account for the required ionizing photons [Muñoz2024]. This issue is further complicated by the need to reconcile various factors such as star formation rates (SFR), escape fractions of ionizing photons, and galaxy properties. Previous research has emphasized the importance of accurately calibrating these parameters to understand reionization [Davies2021], [Park2022], and [Madau2017]. Our work aims to address this photon budget crisis by systematically evaluating the required escape fraction for star-forming galaxies at z\~{}8.
\section{Data and method}
Data and Method:
To tackle this problem, we adopt a literature-anchored approach, relying on established values from prior studies. Specifically, we use the cosmic SFRD analytic fitting function provided by Madau \& Dickinson (2014). For ionizing parameters xi\_ion and O32/beta f\_esc proxy calibrations, we draw from published works such as Chisholm+22, Flury+22, and Simmonds+24. Our method focuses on reconciling the reionization-photon-budget using these literature values without incorporating new observational or catalog data.
\section{Result}
Result:
Our analysis reveals that to close the ionizing-photon-budget at z\~{}8, star-forming galaxies must have an escape fraction of f\_esc=0.419 (+0.422/-0.214). This value is significantly higher than the indirectly inferred f\_esc=0.062 (+0.108/-0.039) derived from LzLCS O32/beta calibrations. The median difference between required and inferred values is +0.330 dex-frac, with 93\% of systematic Monte Carlo simulations showing a shortfall. This result suggests that a genuine shortfall exists in the photon budget, consistent across both O32 and beta calibrations.
\begin{figure}\centering\includegraphics[width=\columnwidth]{result.png}\caption{A residual reionization ionizing-photon-budget shortfall for star-forming galaxies at z\~{}8 (lowxi systematic corner)}\end{figure}
\section{Caveats}
Caveats:
Our study relies on an automated, single-selection approach to reconcile literature values without incorporating new data or calibration. Consequently, our findings are subject to limitations inherent in this method. The systematic uncertainties in xi\_ion, clumping factor (C), and proxy-calibration introduce potential biases that may impact the accuracy of our result. Additionally, the lack of direct observational data and the reliance on published calibrations mean our analysis does not account for variations or updates in these underlying values. Therefore, while our study highlights a significant shortfall in the photon budget, further research incorporating new data and refined calibrations is necessary to confirm and refine this conclusion.
\section*{References}
\footnotesize
[Muoz2024] Muñoz et al. (2024). Reionization after JWST: a photon budget crisis?. bibcode 2024MNRAS.535L..37M \par [Davies2021] Davies et al. (2021). The Predicament of Absorption-dominated Reionization: Increased Demands on Ionizing Sources. bibcode 2021ApJ...918L..35D \par [Park2022] Park et al. (2022). Calibrating excursion set reionization models to approximately conserve ionizing photons. bibcode 2022MNRAS.517..192P \par [Duncan2015] Duncan et al. (2015). Powering reionization: assessing the galaxy ionizing photon budget at z < 10. bibcode 2015MNRAS.451.2030D \par [Madau2017] Madau (2017). Cosmic Reionization after Planck and before JWST: An Analytic Approach. bibcode 2017ApJ...851...50M

\end{document}
